\Problem{Carambola (1)}

\Statement
% de $1.5 \times 3$ metres  

El billar francès o \emph{carambola} és una de les modalitats de jocs de billar existents. En aquest exercici  volem resoldre una versió simplificada d'aquest joc. 
Tenim un jugador i dues boles (blanca i vermella) sobre un tauler entapissat que no té forats a les cantonades. Es col.loca la bola vermella al tauler sense que toqui cap banda, i la bola blanca a una de les quatre  cantonades (vegeu figura).

La bola blanca inicia el recorregut per la diagonal y rebota (seguint l'altre diagonal) quan topa amb una banda.  Si arriba a una cantonada (on nom\'{e}s  hi ha una diagonal) s'atura.

\FigureC{width=8cm}{Rebots1} 

\ \\ 
Per a la simulació assumiren que tenim un tauler de dimensions $n\times m$ i que el recorregut de la  bola blanca acaba quan:

\begin{enumerate}
    \item topa amb una cantonada, 
    \item topa amb la bola vermella,  
    \item fa quatre bandes.
\end{enumerate}

El jugador fa carambola si el recorregut de la bola blanca acaba topant amb la bola vermella després d'haver tocat exactament tres bandes. 

\smallskip
Feu un programa que llegeixi la configuració d'un tauler de billar francès i determini si el jugador fa o no carambola. El programa també ha d'enumerar les posicions on la bola xoca amb les bandes del tauler i mostrar-ne la configuració final. 

%\newpage

%\ \\ 
El vostre programa ha de fer servir les seg\"{u}ents definicions:

\begin{quote}

\textbf{typedef} vector$<$ vector $<$\textbf{char}$>$ $>$ Tauler;

\textbf{struct}  Bola $\{$ 

\hspace{10pt} \textbf{int} x\_ant, y\_ant;\hspace{10pt} // Posicio anterior de la bola

\hspace{10pt} \textbf{int} x\_act, y\_act;\hspace{11pt} // Posicio actual de la bola

\hspace{10pt} \textbf{int} x\_seg, y\_seg;\hspace{10pt} // Posicio seguent de la bola

$\}$;

\end{quote}

%\ \\ 
i tamb\'{e} ha de definir, implementar i utilitzar la funci\'{o}:

\begin{quote}
 $void$ moure\_fins\_xocar ($Tauler$\& t, $Bola$\& b); 
\end{quote}

que, donat un tauler $t$ i la bola blanca $b$,  
mou la bola $b$ dins de $t$ en la direcció i sentit definits per la seva posició actual i següent. La bola s'ha de moure fins a topar amb la bola vermella (si ho fa) o fins arribar a una  banda o cantonada. Quan acaba l'execució de la funció, $b$  ha de tenir com a posició actual  la posició on s'atura la bola, i les posicions prèvia i següent  correctes. 
   
Cal tenir en compte que, quan la bola s'atura a una cantonada, la posició següent és la mateixa posició que l'actual.  A més les posicions recorregudes per la bola han de quedar marcades a $t$. 


%    - xoca_cantonada => xoca_banda
%    - pot ser xoca_bola i xoca_banda (i/o xoca_cantonada) true al mateix temps 

%\newpage

\Input

L'entrada consisteix en diverses descripcions inicials de taulers de billar francès. Cadascuna està formada per dos enters $n, m \geq 2$ que indiquen el nombre de files i col.lumnes de la matriu que representa el tauler, respectivament. A continuació ve una seqüència de quatre enters $x_1, y_1, x_2, y_2$ que descriuen les posicions inicials $(x_1, y_1)$ i $(x_2, y_2)$ de les dues boles. 

\ \\
La bola a la posicio $(x_1, y_1)$ és la bola blanca i sempre està col.locada a una de les cantonades del tauler (és a dir, a les posicions $(0,0)$, 
$(0, m - 1)$, $(n - 1, 0)$ o $(n - 1, m - 1)$). La bola a la posició $(x_2,y_2)$ és la bola vermella i sempre està col.locada a la part interior del tauler, sense tocar les bandes (és a dir, $0<x_2<n-1$ i $0<y_2<m-1$).


\Output

Per a cada tauler d'entrada, escriviu les coordenades de les posicions on la bola blanca topa amb les bandes del tauler. Indiqueu també si fa carambola o no i mostreu-ne la trajectòria descrita. Seguiu el format especificat als exemples. 

%\ \\
\medskip



\Sample
